%% toc: P15 159 159.1,159.2
% Proofs for independence of families of arbitrary size and the complement-choice form.
\section{Families of arbitrary size: proofs}
\label{sec:annex-p15-w02b-p15-006-arbitrary-independence}

Use the definitions of
Chapter~\ref{sec:p15-w02b-p15-006-arbitrary-independence}. The finite
and countable results of
Chapter~\ref{sec:p15-independent-random-variables} are applied only to
families indexed by a finite set, which their hypotheses permit. Two
elementary facts about the symmetric difference are used. For any sets,
\[
 A\triangle C\subseteq(A\triangle B)\cup(B\triangle C):
\]
if \(\omega\in A\setminus C\), then \(\omega\in A\setminus B\) when
\(\omega\notin B\) and \(\omega\in B\setminus C\) when \(\omega\in B\);
the case \(\omega\in C\setminus A\) is the same argument with \(A\) and
\(C\) exchanged. Consequently, for events,
\(\mathbb P(A\triangle C)\leq\mathbb P(A\triangle B)+\mathbb P(B\triangle C)\)
by monotonicity and subadditivity in \cref{thm:p15-measure-laws}.

\begin{annexitem}{annex:p15-w02b-p15-006-finite-character}{thm:p15-w02b-p15-006-finite-character}{Finite subfamilies and countable intersections}
\begin{proof}
In each of the three parts of
\cref{def:p15-w02b-p15-006-arbitrary-families}, the requirement is a
list of equations, one for each finite \(J\subseteq I\) and each
selection of events or value sets at the indices of \(J\). Call these
the \(J\)-equations.

(1) Suppose the family is independent and let \(K\subseteq I\) be
countable. Independence of the subfamily indexed by \(K\), in the sense
of \cref{def:p15-general-independence} or
\cref{def:p15-independent-random-elements}, consists of the
\(J\)-equations for the finite \(J\subseteq K\). Each such \(J\) is a
finite subset of \(I\), so these equations hold. A finite subfamily is
countable, so every finite subfamily is then independent. Conversely,
suppose every finite subfamily is independent, and let \(J\subseteq I\)
be finite. The subfamily indexed by \(J\) is independent, and \(J\) is a
finite subset of its own index set, so the \(J\)-equations hold. As
\(J\) was arbitrary, the family is independent. For a subfamily indexed
by \(J'\subseteq I\), its finite index subsets are finite subsets of
\(I\), so its equations are among those of the whole family.

(2) For countable \(I\), the definitions
\cref{def:p15-general-independence,def:p15-independent-random-elements}
impose exactly the \(J\)-equations for finite \(J\subseteq I\), with the
same conventions for the empty intersection and empty product. The two
definitions therefore impose the same equations.

(3) Put \(E_n=\bigcap_{k=1}^nA_{j_k}\). The finite set
\(\{j_1,\ldots,j_n\}\) has \(n\) distinct elements, so independence gives
\(\mathbb P(E_n)=\prod_{k=1}^n\mathbb P(A_{j_k})\). Passing from \(n\)
to \(n+1\) multiplies by a number in \([0,1]\), so the partial products
do not increase. The events \(E_n\) decrease, and since the enumeration
exhausts \(J\), their intersection is \(\bigcap_{i\in J}A_i\).
Continuity from above for a probability measure,
\cref{thm:p15-measure-laws}, gives
\(\mathbb P(\bigcap_{i\in J}A_i)=\lim_n\mathbb P(E_n)\), which is the
displayed limit. The left side does not mention the enumeration, so
neither does the limit.
\end{proof}
\end{annexitem}

\begin{annexitem}{annex:p15-w02b-p15-006-arbitrary-transfer}{thm:p15-w02b-p15-006-arbitrary-transfer}{Reduction to finite subfamilies}
\begin{proof}
By item~1 of \cref{thm:p15-w02b-p15-006-finite-character}, a family
indexed by \(I\) is independent exactly when the subfamily indexed by
every finite \(J\subseteq I\) is independent. In each item below, fix
such a finite \(J\), apply the finite case of the cited theorem to the
subfamily indexed by \(J\), and then use this equivalence. For
\(J=\varnothing\) every assertion reduces to the empty equation
\(1=1\).

(1) By item~1 of \cref{thm:p15-independence-basic} for the finite family
indexed by \(J\), \((X_i)_{i\in J}\) is independent exactly when
\((\sigma(X_i))_{i\in J}\) is. Since this holds for every finite \(J\),
the whole families are independent together. If the
\(\mathcal F_i\), \(i\in I\), are independent and \(X_i\) is
\(\mathcal F_i\)-measurable, then \((\mathcal F_i)_{i\in J}\) is
independent, and the same finite item makes \((X_i)_{i\in J}\)
independent. If \(\mathcal G_i\subseteq\mathcal F_i\), every finite
selection from the \(\mathcal G_i\) is a selection from the
\(\mathcal F_i\), so its equation holds. Finally, for
\(B\in\mathcal T_i\) the event
\(\{f_i(X_i)\in B\}=\{X_i\in f_i^{-1}(B)\}\) lies in \(\sigma(X_i)\),
because \(f_i^{-1}(B)\in\mathcal S_i\). Thus \(f_i(X_i)\) is
\(\sigma(X_i)\)-measurable, and the \(\sigma(X_i)\) are independent by
the first assertion, so the preceding clause gives independence of the
\(f_i(X_i)\).

(2) The subfamily \((A_i)_{i\in J}\) is a finite independent family,
and \((B_i)_{i\in J}\) is obtained from it by complementing the events
indexed by the finite set \(J\cap K\). Item~2 of
\cref{thm:p15-independence-basic} makes \((B_i)_{i\in J}\) independent.
This holds for every finite \(J\), whether \(K\) is finite or not. The
same item gives, for the finite family indexed by \(J\), the equivalence
of independence of the events, of their indicators, and of the
sigma-algebras \(\{\varnothing,A_i,A_i^c,\Omega\}\). Each of the three
families indexed by \(I\) is independent exactly when all its finite
subfamilies are, so the three conditions are equivalent for \(I\).

(3) For finite \(J\), item~4 of \cref{thm:p15-independence-basic}
applied to \((X_i)_{i\in J}\) and \((Y_i)_{i\in J}\) gives independence
of \((Y_i)_{i\in J}\) and equality of each marginal law. Its proof uses
only the finite union \(\bigcup_{i\in J}N_i\), which is a null event.
The union over all of \(I\) is never formed. In
\cref{def:p15-w02b-p15-006-mixed-family} the exceptional sets are the
singletons \(\{t\}\), \(t\in(0,1)\), and their union is \((0,1)\), of
probability one.

(4) The collections \(\mathcal C_i\), \(i\in J\), are finitely many
independent pi-systems. \Cref{thm:p15-independence-generators} makes
\((\sigma(\mathcal C_i))_{i\in J}\) independent.

(5) Suppose the \(X_i\) are independent and \(i_1,\ldots,i_n\) are
distinct. The subfamily indexed by \(\{i_1,\ldots,i_n\}\) is
independent, and the necessity part of
\cref{thm:p15-independence-thresholds} gives the threshold equations,
since rational thresholds are among the real ones. Conversely, assume
the rational threshold equations for every finite list of distinct
indices. For a nonempty finite \(J\), list its elements without
repetition; the sufficiency part of
\cref{thm:p15-independence-thresholds}, with rational thresholds, makes
\((X_i)_{i\in J}\) independent.

(6) The last assertion of \cref{thm:p15-independent-expectations},
applied to the finite family \((X_i)_{i\in J}\), states that this
family is independent exactly when the bounded-function product
equation holds on every subset \(J'\subseteq J\). If the whole family is
independent, its subfamily indexed by \(J\) is, so the equation holds on
\(J'=J\). Conversely, if the equation holds for every finite subset of
\(I\), it holds for every \(J'\subseteq J\), so \((X_i)_{i\in J}\) is
independent.

(7) Fix a nonempty finite \(J\subseteq I\) and list it without
repetition as \(i_1,\ldots,i_n\). The countable-value clause of
\cref{thm:p15-independent-product-law}, applied to
\(X_{i_1},\ldots,X_{i_n}\), states that this finite family is
independent exactly when the point-mass equation holds for all
\(x_k\in S_{i_k}\). Since the whole family is independent exactly when
each such finite subfamily is, the point-mass equations for all finite
lists of distinct indices are equivalent to independence of
\((X_i)_{i\in I}\). A list of distinct indices in another order gives
the same equation with its factors permuted, so the choice of listing
does not matter.
\end{proof}
\end{annexitem}

\begin{annexitem}{annex:p15-w02c-p15-006-lambda-class}{thm:p15-w02c-p15-006-lambda-class}{The factorization class}
\begin{proof}
(1) Since \(\Omega\cap D=D\) and \(\mathbb P(\Omega)=1\),
\(\mathbb P(\Omega\cap D)=\mathbb P(\Omega)\mathbb P(D)\), so
\(\Omega\in\mathcal L_D\). Let \(A\subseteq B\) lie in
\(\mathcal L_D\). Then \(A\cap D\subseteq B\cap D\) and
\((B\setminus A)\cap D=(B\cap D)\setminus(A\cap D)\). Item~1 of
\cref{thm:p15-measure-laws}, applied twice to finite measures, gives
\[
 \mathbb P((B\setminus A)\cap D)
   =\mathbb P(B\cap D)-\mathbb P(A\cap D)
   =(\mathbb P(B)-\mathbb P(A))\mathbb P(D)
   =\mathbb P(B\setminus A)\mathbb P(D),
\]
so \(B\setminus A\in\mathcal L_D\). Let
\(B_1\subseteq B_2\subseteq\cdots\) lie in \(\mathcal L_D\) and put
\(B=\bigcup_rB_r\). The events \(B_r\cap D\) increase, with union
\(B\cap D\). Item~3 of \cref{thm:p15-measure-laws}, applied to both
sequences, gives
\[
 \mathbb P(B\cap D)=\lim_r\mathbb P(B_r\cap D)
     =\lim_r\mathbb P(B_r)\mathbb P(D)=\mathbb P(B)\mathbb P(D),
\]
so \(B\in\mathcal L_D\). These are the three requirements of
\cref{def:p15-pi-lambda}. For \(A\in\mathcal L_D\), the difference
clause with \(A\subseteq\Omega\) gives \(A^c=\Omega\setminus A\in\mathcal L_D\).

(2) In the four-point model every outcome has probability \(1/4\), so
\(\mathbb P(C)=1/2\). Since \(A\cap C=\{1\}\) and
\(\mathbb P(A)=1/2\),
\[
 \mathbb P(A\cap C)=\tfrac14=\tfrac12\cdot\tfrac12=\mathbb P(A)\mathbb P(C),
\]
and the same computation with \(B\cap C=\{1\}\) puts \(B\) in
\(\mathcal L_C\). For the intersection,
\((A\cap B)\cap C=\{1\}\) has probability \(1/4\), whereas
\(\mathbb P(A\cap B)\mathbb P(C)=\tfrac14\cdot\tfrac12=\tfrac18\). For
the union, \((A\cup B)\cap C=\{1\}\) has probability \(1/4\), whereas
\(\mathbb P(A\cup B)\mathbb P(C)=\tfrac34\cdot\tfrac12=\tfrac38\). Hence
neither \(A\cap B\) nor \(A\cup B\) lies in \(\mathcal L_C\). A
sigma-algebra contains the intersection of any two of its members, so
\(\mathcal L_C\) is not a sigma-algebra.

(3) The sigma-algebra \(\sigma(\{A,B\})\) contains \(A\) and \(B\) and
hence \(A\cap B\), which by item~2 is not in \(\mathcal L_C\). So
\(\{A,B\}\subseteq\mathcal L_C\) while
\(\sigma(\{A,B\})\not\subseteq\mathcal L_C\). If instead a pi-system
\(\mathcal C\) satisfies \(\mathcal C\subseteq\mathcal L_D\), then
\(\mathcal L_D\) is a lambda-system containing \(\mathcal C\) by
item~1, and \cref{thm:p15-measure-uniqueness} gives
\(\sigma(\mathcal C)\subseteq\mathcal L_D\).
\end{proof}
\end{annexitem}

\begin{annexitem}{annex:p15-w02b-p15-006-arbitrary-blocks}{thm:p15-w02b-p15-006-arbitrary-blocks}{Finite intersections inside groups}
\begin{proof}
For \(j\in L\), let \(\mathcal C_j\) consist of all finite intersections
\(\bigcap_{i\in K}A_i\) with \(K\subseteq I_j\) finite and
\(A_i\in\mathcal F_i\) for \(i\in K\); the empty intersection
\(\Omega\) is included. It is a pi-system: given
\(\bigcap_{i\in K}A_i\) and \(\bigcap_{i\in K'}A'_i\), their
intersection is \(\bigcap_{i\in K\cup K'}A''_i\), where
\(A''_i=A_i\cap A'_i\) for \(i\in K\cap K'\), \(A''_i=A_i\) for
\(i\in K\setminus K'\), and \(A''_i=A'_i\) for \(i\in K'\setminus K\).
Each \(A''_i\) lies in the sigma-algebra \(\mathcal F_i\), and
\(K\cup K'\) is a finite subset of \(I_j\).

Taking one-element \(K\) shows
\(\bigcup_{i\in I_j}\mathcal F_i\subseteq\mathcal C_j\). Every member of
\(\mathcal C_j\) is a finite intersection of members of that union and
so lies in \(\mathcal G_j\). Hence
\(\mathcal G_j\subseteq\sigma(\mathcal C_j)\subseteq\mathcal G_j\),
and \(\sigma(\mathcal C_j)=\mathcal G_j\). When \(I_j\) is empty,
\(\mathcal C_j=\{\Omega\}\) and \(\sigma(\mathcal C_j)=\{\varnothing,\Omega\}\),
in agreement with the stated convention.

The collections \(\mathcal C_j\), \(j\in L\), are independent. Let
\(L'\subseteq L\) be finite and choose \(C_j\in\mathcal C_j\) for
\(j\in L'\), each with a representation
\(C_j=\bigcap_{i\in K_j}A_i\), \(K_j\subseteq I_j\) finite,
\(A_i\in\mathcal F_i\). The groups are disjoint, so the \(K_j\) are
disjoint and \(K=\bigcup_{j\in L'}K_j\) is a finite subset of \(I\)
whose elements each carry exactly one event \(A_i\). Independence of
the \(\mathcal F_i\), first over \(K\) and then over each \(K_j\), gives
\[
 \mathbb P\left(\bigcap_{j\in L'}C_j\right)
   =\mathbb P\left(\bigcap_{i\in K}A_i\right)
   =\prod_{j\in L'}\prod_{i\in K_j}\mathbb P(A_i)
   =\prod_{j\in L'}\mathbb P(C_j).
\]
Empty \(K_j\) contribute the factor one on both sides. Item~4 of
\cref{thm:p15-w02b-p15-006-arbitrary-transfer}, which allows an
arbitrary index set \(L\), now makes the \(\sigma(\mathcal C_j)=\mathcal G_j\)
independent. If each \(Y_j\) is \(\mathcal G_j\)-measurable, item~1 of
the same theorem gives independence of the \(Y_j\).
\end{proof}
\end{annexitem}

\begin{annexitem}{annex:p15-w02b-p15-006-uniform-approximation}{thm:p15-w02b-p15-006-uniform-approximation}{Approximation by cells}
\begin{proof}
For a fixed \(m\) there are \(2^m\) subsets \(K\), so the sets of mesh
\(1/m\) form a finite collection \(\mathcal D_m\), and
\(\mathcal D=\bigcup_{m\geq1}\mathcal D_m\) is a countable union of
finite collections indexed by \(m\); listing \(\mathcal D_1\), then
\(\mathcal D_2\), and so on enumerates it. Each cell is a Borel
interval, a finite union of cells is Borel, and its intersection with
\((0,1)\) belongs to \(\mathcal F_U\) by the definition of
\(\mathcal F_U\) in \cref{thm:p15-lebesgue-unit-interval}.

Three closure properties of \(\mathcal D\) are needed. First, the cells
of mesh \(1/m\) are pairwise disjoint with union \((0,1]\): a point
\(u\in(0,1]\) lies in the cell with index \(k\) exactly when
\(k<mu\leq k+1\), and exactly one integer \(k\in\{0,\ldots,m-1\}\)
satisfies this. Hence
\[
 (0,1)\setminus\Bigl((0,1)\cap\bigcup_{k\in K}(k/m,(k+1)/m]\Bigr)
   =(0,1)\cap\bigcup_{k\notin K}(k/m,(k+1)/m],
\]
so \(\mathcal D\) is closed under complements relative to \((0,1)\).
Second, for integers \(m,m'\geq1\), each cell of mesh \(1/m\) is the
union of the \(m'\) cells
\(((km'+r)/(mm'),(km'+r+1)/(mm')]\), \(0\leq r<m'\), of mesh
\(1/(mm')\). Two members of \(\mathcal D\), of meshes \(1/m\) and
\(1/m'\), are therefore both of mesh \(1/(mm')\), and their union is the
member of that mesh with the union of the two index sets. By induction
on the number of sets, \(\mathcal D\) is closed under finite unions.
Third, for rational \(q\), the set \((-\infty,q]\cap(0,1)\) belongs to
\(\mathcal D\): it is empty, the member with \(K=\varnothing\), when
\(q\leq0\); it is \((0,1)\), the member with \(m=1\) and \(K=\{0\}\),
when \(q\geq1\); and when \(0<q<1\), write \(q=r/s\) with integers
\(1\leq r<s\), so that it equals
\((0,r/s]=(0,1)\cap\bigcup_{k=0}^{r-1}(k/s,(k+1)/s]\).

Let \(\mathcal G\) be the class of Borel sets \(B\subseteq\mathbb R\)
such that for every \(\varepsilon>0\) some \(D\in\mathcal D\) satisfies
\(\mathbb U((B\cap(0,1))\triangle D)<\varepsilon\). We show it is a
sigma-algebra.
\begin{itemize}
\item \(\mathbb R\in\mathcal G\), with \(D=(0,1)\) and empty symmetric
      difference.
\item Let \(B\in\mathcal G\), \(A=B\cap(0,1)\), and let \(D\in\mathcal D\)
      satisfy \(\mathbb U(A\triangle D)<\varepsilon\). Then
      \((\mathbb R\setminus B)\cap(0,1)=(0,1)\setminus A\), and
      \(D'=(0,1)\setminus D\in\mathcal D\) by the first closure property.
      Since \(A\) and \(D\) are subsets of \((0,1)\), a point of
      \((0,1)\) lies in \(((0,1)\setminus A)\setminus((0,1)\setminus D)\)
      exactly when it lies in \(D\setminus A\), and symmetrically; points
      outside \((0,1)\) lie in neither set. Thus
      \(((0,1)\setminus A)\triangle D'=A\triangle D\), and
      \(\mathbb R\setminus B\in\mathcal G\).
\item Let \(B_n\in\mathcal G\), \(n\geq1\), put \(A_n=B_n\cap(0,1)\),
      \(A=\bigcup_nA_n=(\bigcup_nB_n)\cap(0,1)\), and
      \(H_N=\bigcup_{n\leq N}A_n\). The events \(H_N\) increase to
      \(A\), so by continuity from below and finiteness of
      \(\mathbb U\) there is \(N\) with
      \(\mathbb U(A\setminus H_N)=\mathbb U(A)-\mathbb U(H_N)<\varepsilon/2\).
      For \(n\leq N\) choose \(D_n\in\mathcal D\) with
      \(\mathbb U(A_n\triangle D_n)<\varepsilon/(2N)\), and put
      \(D=\bigcup_{n\leq N}D_n\), a member of \(\mathcal D\) by the
      second closure property. Then
      \[
       A\triangle D\subseteq(A\setminus H_N)\cup
                     \bigcup_{n\leq N}(A_n\triangle D_n).
      \]
      Indeed, a point of \(A\setminus D\) either lies outside \(H_N\)
      or lies in some \(A_n\), \(n\leq N\), but not in \(D_n\); a point of
      \(D\setminus A\) lies in some \(D_n\), \(n\leq N\), but not in
      \(A_n\). Subadditivity gives
      \(\mathbb U(A\triangle D)<\varepsilon/2+N\varepsilon/(2N)=\varepsilon\),
      so \(\bigcup_nB_n\in\mathcal G\).
\end{itemize}
By the third closure property, \(\mathcal G\) contains every ray
\((-\infty,q]\), \(q\in\mathbb Q\), with zero error. These rays generate
\(\mathcal B(\mathbb R)\) by \cref{thm:p15-borel-generators}, so
\(\mathcal G=\mathcal B(\mathbb R)\). Every \(A\in\mathcal F_U\) has the
form \(B\cap(0,1)\) with \(B\) Borel, which proves the approximation.
\end{proof}
\end{annexitem}

\begin{annexitem}{annex:p15-w02b-p15-006-separable-limit}{thm:p15-w02b-p15-006-separable-limit}{Separated events and an injective code}
\begin{proof}
The collection \(\mathcal D\) is nonempty, since \(\Omega\) must be
approximated. Fix an enumeration \(d_1,d_2,\ldots\) of \(\mathcal D\),
finite or infinite.

(1) Put \(p=\mathbb P(A)\) and \(q=\mathbb P(B)\). The sets
\(A\setminus B=A\cap B^c\) and \(B\setminus A=A^c\cap B\) are disjoint.
For an independent pair, \(A,B^c\) and \(A^c,B\) are independent pairs by
\cref{thm:p15-finite-independence-laws} (equivalently item~2 of
\cref{thm:p15-independence-basic}). Hence
\[
 \mathbb P(A\triangle B)=p(1-q)+(1-p)q\geq\delta\cdot\delta+\delta\cdot\delta
                       =2\delta^2,
\]
because each of \(p,1-p,q,1-q\) is at least \(\delta\).

(2) For \(i\in I\), let \(\varphi(i)\) be the least \(r\) with
\(\mathbb P(A_i\triangle d_r)<\delta^2\); it exists by the approximation
hypothesis. Suppose \(i\neq k\) and \(\varphi(i)=\varphi(k)=r\). The
symmetric-difference inequality gives
\(\mathbb P(A_i\triangle A_k)\leq\mathbb P(A_i\triangle d_r)
+\mathbb P(d_r\triangle A_k)<2\delta^2\). But \(A_i,A_k\) form an
independent pair, a two-element subfamily, with probabilities in
\([\delta,1-\delta]\), so part~(1) gives
\(\mathbb P(A_i\triangle A_k)\geq2\delta^2\). This contradiction shows
that \(\varphi\) is injective from \(I\) into the positive integers,
so \(I\) is countable.

(3) Let \(I^*\) be the set of indices with \(0<\mathbb P(A_i)<1\). For
\(i\in I^*\) let \(m(i)\) be the least integer \(m\geq2\) with
\(1/m\leq\mathbb P(A_i)\leq1-1/m\); it exists because any
\(m\geq\max\{2,1/\mathbb P(A_i),1/(1-\mathbb P(A_i))\}\) qualifies. Let
\(\varphi(i)\) be the least \(r\) with
\(\mathbb P(A_i\triangle d_r)<1/m(i)^2\). If \(i\neq k\) in \(I^*\)
had \(m(i)=m(k)=m\) and \(\varphi(i)=\varphi(k)\), the argument of
part~(2) with \(\delta=1/m\leq1/2\) would give a contradiction. Thus
\(i\mapsto(m(i),\varphi(i))\) is injective from \(I^*\) into the
countable set of pairs of positive integers.

(4) First the criterion for constancy. If \(X=c\) outside a null event
\(N\), then \(\{X\leq q\}\) agrees outside \(N\) with \(\Omega\) when
\(c\leq q\) and with \(\varnothing\) when \(c>q\), so
\(\mathbb P(X\leq q)\in\{0,1\}\). Conversely assume
\(\mathbb P(X\leq q)\in\{0,1\}\) for all rational \(q\). The events
\(\{X\leq n\}\) increase to \(\Omega\) and \(\{X\leq-n\}\) decrease to
\(\varnothing\), as \(X\) is real-valued, so by continuity some integer
\(n\) has \(\mathbb P(X\leq n)>0\), hence equal to one, and some integer
\(n'\) has \(\mathbb P(X\leq-n')<1\), hence equal to zero. The set
\(Q_1=\{q\in\mathbb Q:\mathbb P(X\leq q)=1\}\) is therefore nonempty, and
by monotonicity it contains no \(q\leq-n'\). Its infimum \(c\) is real.
For rational \(q>c\), some \(q'\in Q_1\) satisfies \(q'<q\), and
\(\mathbb P(X\leq q)\geq\mathbb P(X\leq q')=1\). For rational \(q<c\),
\(q\notin Q_1\), so \(\mathbb P(X\leq q)=0\). Choose rationals
\(q_r\downarrow c\) with \(q_r>c\) and rationals \(p_r\uparrow c\) with
\(p_r<c\), both strictly monotone. Then \(\{X\leq q_r\}\) decreases to
\(\{X\leq c\}\) and \(\{X\leq p_r\}\) increases to \(\{X<c\}\): a value
\(X(\omega)<c\) is at most \(p_r\) once \(p_r\) is close enough to
\(c\). Continuity from above and below gives \(\mathbb P(X\leq c)=1\)
and \(\mathbb P(X<c)=0\), hence \(\mathbb P(X=c)=1\).

Now let \((X_i)_{i\in I}\) be independent and let \(I^*\) be the set of
indices at which \(X_i\) is not almost surely constant. Fix an
enumeration \(q_1,q_2,\ldots\) of \(\mathbb Q\). For \(i\in I^*\), the
criterion gives a least \(s(i)\) with
\(0<\mathbb P(X_i\leq q_{s(i)})<1\). Put
\(E_i=\{X_i\leq q_{s(i)}\}\), and define \(m(i)\) and \(\varphi(i)\)
from \(E_i\) as in part~(3). If \(i\neq k\) in \(I^*\) shared the
triple \((s,m,\varphi)\), then \(E_i\in\sigma(X_i)\) and
\(E_k\in\sigma(X_k)\) would be independent events by item~1 of
\cref{thm:p15-w02b-p15-006-arbitrary-transfer}, with probabilities in
\([1/m,1-1/m]\) and \(\mathbb P(E_i\triangle E_k)<2/m^2\), contradicting
part~(1). So \(i\mapsto(s(i),m(i),\varphi(i))\) is injective into a
countable set of triples of positive integers, and \(I^*\) is
countable.
\end{proof}
\end{annexitem}

\begin{annexitem}{annex:p15-w02b-p15-006-uniform-countable}{thm:p15-w02b-p15-006-uniform-countable}{The uniform space}
\begin{proof}
\Cref{thm:p15-w02b-p15-006-uniform-approximation} supplies a countable
collection \(\mathcal D\subseteq\mathcal F_U\) that approximates every
event of \(\mathcal F_U\) in \(\mathbb U\)-measure of symmetric
difference. Parts~(3) and~(4) of
\cref{thm:p15-w02b-p15-006-separable-limit} therefore apply to
\((\Omega_U,\mathcal F_U,\mathbb U)\) and give the two countability
assertions. If a family on some probability space is independent and
has uncountably many members of probability strictly between zero and
one, or uncountably many real members that are not almost surely
constant, the same two parts show that the space cannot have a
countable approximating collection.
\end{proof}
\end{annexitem}

\begin{annexitem}{annex:p15-w02b-p15-006-complement-form}{thm:p15-w02b-p15-006-complement-form}{Complement choices}
\begin{proof}
For an event \(A\) write \(A^{(1)}=A\) and \(A^{(0)}=A^c\).

(1)\(\Rightarrow\)(3). Assume the subset form and let
\(C_j\in\{A_j,A_j^c,\Omega\}\). Put \(M=\{j:C_j=\Omega\}\), let
\(B_j=C_j\) for \(j\notin M\) and \(B_j=A_j\) for \(j\in M\). The
events \(B_1,\ldots,B_n\) arise from the independent events
\(A_1,\ldots,A_n\) by complementing those with \(C_j=A_j^c\), so they
are independent by item~2 of \cref{thm:p15-independence-basic}. The
subset equation for \(\{1,\ldots,n\}\setminus M\) and the identity
\(\mathbb P(\Omega)=1\) give
\[
 \mathbb P\left(\bigcap_{j=1}^nC_j\right)
   =\mathbb P\left(\bigcap_{j\notin M}B_j\right)
   =\prod_{j\notin M}\mathbb P(B_j)
   =\prod_{j=1}^n\mathbb P(C_j).
\]
This is the argument of Durrett's Theorem~2.1.3
\cite[pp.~44--45]{DurrettPTEAuthor2019}.

(3)\(\Rightarrow\)(2). The choices without \(\Omega\) are among those
of the extended form.

(2)\(\Rightarrow\)(1). Fix \(J\subseteq\{1,\ldots,n\}\) and let
\(M=\{1,\ldots,n\}\setminus J\). For each function
\(\varepsilon:M\to\{0,1\}\) put
\[
 E_\varepsilon=\bigcap_{j\in J}A_j\cap\bigcap_{k\in M}A_k^{(\varepsilon(k))}.
\]
These \(2^{\#M}\) events are pairwise disjoint: if
\(\varepsilon(k)\neq\varepsilon'(k)\), then \(E_\varepsilon\subseteq
A_k^{(\varepsilon(k))}\) and \(E_{\varepsilon'}\subseteq
A_k^{(\varepsilon'(k))}\), which are complementary. Their union is
\(\bigcap_{j\in J}A_j\): each \(E_\varepsilon\) is contained in it, and a
point \(\omega\) of it lies in \(E_\varepsilon\) for the function with
\(\varepsilon(k)=1\) exactly when \(\omega\in A_k\). Each
\(E_\varepsilon\) is a full intersection of a complement choice, with
\(B_j=A_j\) for \(j\in J\) and \(B_k=A_k^{(\varepsilon(k))}\) for
\(k\in M\). Finite additivity and the complement-choice form give
\[
 \mathbb P\left(\bigcap_{j\in J}A_j\right)
   =\sum_{\varepsilon}\mathbb P(E_\varepsilon)
   =\prod_{j\in J}\mathbb P(A_j)\,
    \sum_{\varepsilon}\prod_{k\in M}\mathbb P\bigl(A_k^{(\varepsilon(k))}\bigr).
\]
For real numbers \(a_k,b_k\), \(k\in M\),
\[
 \prod_{k\in M}(a_k+b_k)=\sum_{\varepsilon:M\to\{0,1\}}
       \prod_{k\in M}c_k(\varepsilon(k)),
 \qquad c_k(1)=a_k,\ c_k(0)=b_k.
\]
This holds for \(M=\varnothing\), where both sides are one, and passes
from \(M\) to \(M\cup\{k_0\}\), \(k_0\notin M\), on multiplying by
\(a_{k_0}+b_{k_0}\) and splitting each function on \(M\cup\{k_0\}\)
according to its value at \(k_0\). With \(a_k=\mathbb P(A_k)\) and
\(b_k=\mathbb P(A_k^c)\), every factor \(a_k+b_k\) equals one, so the
sum over \(\varepsilon\) is one and
\(\mathbb P(\bigcap_{j\in J}A_j)=\prod_{j\in J}\mathbb P(A_j)\). The
cases \(J=\varnothing\) and \(M=\varnothing\) are included.

For an arbitrary family \((A_i)_{i\in I}\), item~1 of
\cref{thm:p15-w02b-p15-006-finite-character} makes independence
equivalent to independence of every finite subfamily. For a nonempty
finite \(J\subseteq I\), list its elements without repetition; the
equivalence of~(1) and~(2) just proved shows that the subfamily
indexed by \(J\) is independent exactly when its complement-choice
equations hold. For \(J=\varnothing\) both conditions are the equation
\(1=1\).
\end{proof}
\end{annexitem}
